% Options for packages loaded elsewhere
\PassOptionsToPackage{unicode}{hyperref}
\PassOptionsToPackage{hyphens}{url}
\documentclass[
  11pt,
]{article}
\usepackage{xcolor}
\usepackage[left=3cm,right=3cm,top=2.5cm,bottom=2.5cm]{geometry}
\usepackage{amsmath,amssymb}
\setcounter{secnumdepth}{5}
\usepackage{iftex}
\ifPDFTeX
  \usepackage[T1]{fontenc}
  \usepackage[utf8]{inputenc}
  \usepackage{textcomp} % provide euro and other symbols
\else % if luatex or xetex
  \usepackage{unicode-math} % this also loads fontspec
  \defaultfontfeatures{Scale=MatchLowercase}
  \defaultfontfeatures[\rmfamily]{Ligatures=TeX,Scale=1}
\fi
\usepackage{lmodern}
\ifPDFTeX\else
  % xetex/luatex font selection
\fi
% Use upquote if available, for straight quotes in verbatim environments
\IfFileExists{upquote.sty}{\usepackage{upquote}}{}
\IfFileExists{microtype.sty}{% use microtype if available
  \usepackage[]{microtype}
  \UseMicrotypeSet[protrusion]{basicmath} % disable protrusion for tt fonts
}{}
\usepackage{setspace}
\makeatletter
\@ifundefined{KOMAClassName}{% if non-KOMA class
  \IfFileExists{parskip.sty}{%
    \usepackage{parskip}
  }{% else
    \setlength{\parindent}{0pt}
    \setlength{\parskip}{6pt plus 2pt minus 1pt}}
}{% if KOMA class
  \KOMAoptions{parskip=half}}
\makeatother
\usepackage{longtable,booktabs,array}
\usepackage{caption}
\captionsetup[table]{skip=6pt}
\newcounter{none} % for unnumbered tables
\usepackage{calc} % for calculating minipage widths
% Correct order of tables after \paragraph or \subparagraph
\usepackage{etoolbox}
\makeatletter
\patchcmd\longtable{\par}{\if@noskipsec\mbox{}\fi\par}{}{}
\makeatother
% Allow footnotes in longtable head/foot
\IfFileExists{footnotehyper.sty}{\usepackage{footnotehyper}}{\usepackage{footnote}}
\makesavenoteenv{longtable}
\usepackage{graphicx}
\makeatletter
\newsavebox\pandoc@box
\newcommand*\pandocbounded[1]{% scales image to fit in text height/width
  \sbox\pandoc@box{#1}%
  \Gscale@div\@tempa{\textheight}{\dimexpr\ht\pandoc@box+\dp\pandoc@box\relax}%
  \Gscale@div\@tempb{\linewidth}{\wd\pandoc@box}%
  \ifdim\@tempb\p@<\@tempa\p@\let\@tempa\@tempb\fi% select the smaller of both
  \ifdim\@tempa\p@<\p@\scalebox{\@tempa}{\usebox\pandoc@box}%
  \else\usebox{\pandoc@box}%
  \fi%
}
% Set default figure placement to htbp
\def\fps@figure{htbp}
\makeatother
\setlength{\emergencystretch}{3em} % prevent overfull lines
\providecommand{\tightlist}{%
  \setlength{\itemsep}{0pt}\setlength{\parskip}{0pt}}
%% tools/stamp.tex
%% Provenance footer for PDFs rendered from this project.
%% Vendored by zzcollab; this file is not project-specific. The
%% footer prints the source document and the time of rendering.
%% The git version is defined here too, and so is available to a
%% project that wants it on the page, but it is left off the
%% default footer: it already names the staged copy in share/ and
%% appears in its manifest row. All values are supplied at render
%% time by tools/stamp-render.R, which writes a small generated
%% file that redefines the macros below. The \providecommand
%% defaults keep a bare LaTeX compile working when the document is
%% built without the wrapper.
\usepackage{fancyhdr}
\usepackage{xcolor}
\providecommand{\stampsource}{\detokenize{(source not recorded)}}
\providecommand{\stamptime}{(time not recorded)}
\providecommand{\stampversion}{\detokenize{(version not recorded)}}
\setlength{\footskip}{40pt}
\newcommand{\stampfootcontent}{%
  \footnotesize\color{gray}%
  \parbox[b]{\textwidth}{%
    \stamptime\hfill\thepage\\[2pt]%
    \ttfamily\raggedright\stampsource}}
\pagestyle{fancy}
\fancyhf{}
\fancyfoot[C]{\stampfootcontent}
\renewcommand{\headrulewidth}{0pt}
\renewcommand{\footrulewidth}{0.4pt}
%% The title page uses the `plain` page style; redefine it so the
%% stamp appears there too.
\fancypagestyle{plain}{%
  \fancyhf{}%
  \fancyfoot[C]{\stampfootcontent}%
  \renewcommand{\headrulewidth}{0pt}%
  \renewcommand{\footrulewidth}{0.4pt}}
\renewcommand{\stampsource}{\textasciitilde{}/\allowbreak{}prj/\allowbreak{}res/\allowbreak{}36-pmsimstats-ng/\allowbreak{}pmsimstats-ng/\allowbreak{}analysis/\allowbreak{}report/\allowbreak{}11-combined-dgp-architecture/\allowbreak{}bullets.Rmd}
\renewcommand{\stamptime}{2026-08-15 16:07 PDT}
\renewcommand{\stampversion}{\detokenize{bc0b671-wip}}
\usepackage{amsmath}
\usepackage{amssymb}
\usepackage{booktabs}
\usepackage{longtable}
\usepackage{float}
\usepackage[mathlines]{lineno}
\linenumbers
\usepackage{bookmark}
\IfFileExists{xurl.sty}{\usepackage{xurl}}{} % add URL line breaks if available
\urlstyle{same}
% fallback for those not using the hyperref driver hyperxmp:
\makeatletter
\@ifundefined{xmpquote}{\newcommand{\xmpquote}[1]{#1}}{}
\makeatother
\hypersetup{
  pdftitle={Structured summary: A combined data-generating architecture for biomarker-treatment interactions: channel super-additivity and mixing-weight sensitivity in aggregated N-of-1 trials},
  pdfauthor={pmsimstats team},
  hidelinks,
  pdfcreator={LaTeX via pandoc}}

\title{Structured summary: A combined data-generating architecture for biomarker-treatment interactions: channel super-additivity and mixing-weight sensitivity in aggregated N-of-1 trials}
\author{pmsimstats team}
\date{August 15, 2026}

\begin{document}
\maketitle

{
\setcounter{tocdepth}{2}
\tableofcontents
}
\setstretch{1.1}
\emph{This document is a structured, section-by-section summary of the key points argued in the companion manuscript, ``A combined data-generating architecture for biomarker-treatment interactions: channel super-additivity and mixing-weight sensitivity in aggregated N-of-1 trials.'' It is provided as a supplementary reading aid and does not stand on its own; all notation, derivations, simulation detail, and citations are in the main manuscript.}

\bigskip\noindent

\rule{\textwidth}{0.4pt}\bigskip

\section{1. Introduction}

\textbf{The companion paper contrasts two data-generating process
(DGP) architectures; this paper treats the case where both operate
simultaneously.}

\begin{itemize}\setlength{\itemsep}{2pt}\setlength{\parskip}{0pt}
\item The mean-moderation architecture encodes the interaction in the mean;
  the covariance-moderation architecture encodes it in the covariance.
\item A biomarker could act through both pathways at once.
\item The dual-channel architecture carries independent weights on the two channels
  and recovers A and B as its single-channel limits.
\end{itemize}

\textbf{The paper makes two contributions: an analytical
super-additivity result and a calibrated simulation of the combined
grid.}

\begin{itemize}\setlength{\itemsep}{2pt}\setlength{\parskip}{0pt}
\item Analytical: the channel signals add, and convexity of the
  power function makes their power combination super-additive.
\item Empirical: a $3 \times 3$ channel-weight grid crossed with
  design and carryover, with boundary cells validated against the
  companion paper's pure-architecture baselines.
\item Guidance: when a combined model is warranted and how to use it
  as a sensitivity framework.
\end{itemize}

\subsection{2.1 Definition}

\textbf{The dual-channel architecture applies both interaction mechanisms in
sequence: a differential-correlation multivariate normal (MVN) draw,
then an on-drug mean shift.}

\begin{itemize}\setlength{\itemsep}{2pt}\setlength{\parskip}{0pt}
\item The MVN covariance encodes BM-BR differential correlation with
  weight $c_{bm,B}$ (the covariance-moderation architecture channel).
\item After the draw, BR is additively shifted on-drug by
  $c_{bm,A} \cdot z_B \cdot \sigma_{BR}$ (the mean-moderation architecture
  channel).
\item The two weights are independent and are calibrated separately.
\end{itemize}

\subsection{2.2 Boundary reductions and identifiability}

\textbf{The two channels are algebraically independent, so the
boundary cases recover the pure architectures exactly.}

\begin{itemize}\setlength{\itemsep}{2pt}\setlength{\parskip}{0pt}
\item $c_{bm,B} = 0$ removes the covariance channel and recovers
  the pure mean-moderation architecture.
\item $c_{bm,A} = 0$ removes the mean-shift channel and recovers
  the pure covariance-moderation architecture.
\item The boundaries are enforced by construction, so they serve as
  internal validation gates.
\end{itemize}

\subsection{2.3 Biological rationale}

\textbf{The two channels are mechanistically distinct and need not
scale in proportion.}

\begin{itemize}\setlength{\itemsep}{2pt}\setlength{\parskip}{0pt}
\item The mean-shift channel (A) reflects pharmacodynamic scaling:
  the drug's effect is larger in high-biomarker participants.
\item The covariance channel (B) reflects differential
  co-regulation that couples biomarker and response on-drug and
  washes out off-drug.
\item Neither channel implies the other; a drug could activate one,
  both, or neither.
\end{itemize}

\section{3. Channel Super-additivity: an Analytical Account}

\textbf{The biomarker-treatment signal delivered to the analysis
model is additive across the two channels.}

\begin{itemize}\setlength{\itemsep}{2pt}\setlength{\parskip}{0pt}
\item The interaction coefficient depends on the on-drug covariance
  $\mathrm{Cov}(B, BR)$.
\item The covariance channel contributes
  $c_{bm,B}\,\sigma_B \sigma_{BR}$; the mean channel contributes
  $c_{bm,A}\,\sigma_B \sigma_{BR}$.
\item The total on-drug signal covariance is
  $(c_{bm,A} + c_{bm,B})\,\sigma_B \sigma_{BR}$: the channels add.
\end{itemize}

\textbf{Additive signal becomes super-additive power because the
power function is convex over the sub-ceiling regime.}

\begin{itemize}\setlength{\itemsep}{2pt}\setlength{\parskip}{0pt}
\item Power is a convex function of the noncentrality where power is
  below roughly one half.
\item Convexity gives
  $\pi(\eta_A + \eta_B) - \pi(0) >
   [\pi(\eta_A) - \pi(0)] + [\pi(\eta_B) - \pi(0)]$.
\item Hence interior cells exceed the probability-scale additive
  benchmark; the effect attenuates and can reverse near the ceiling.
\end{itemize}

\subsection{4.1 Design}

\textbf{A $3 \times 3$ channel-weight grid crossed with three
designs and three carryover half-lives.}

\begin{itemize}\setlength{\itemsep}{2pt}\setlength{\parskip}{0pt}
\item Channel weights: $c_{bm,A} \times c_{bm,B} \in
  \{0, 0.22, 0.45\}^2$.
\item Designs: CO, Hybrid, OL+BDC; carryover
  $t_{1/2} \in \{0, 0.5, 1.0\}$ weeks.
\item E2 analysis specification; 1000 replicates per cell.
\end{itemize}

\subsection{4.2 Boundary validation against the companion baselines}

\textbf{The boundary cells are validated numerically against the
companion paper's pure-architecture power values.}

\begin{itemize}\setlength{\itemsep}{2pt}\setlength{\parskip}{0pt}
\item The combined driver is run under the same sample-size and
  calibration settings as the companion paper's Section 3.1 driver.
\item The $c_{bm,B} = 0$ column must reproduce the pure mean-moderation architecture;
  the $c_{bm,A} = 0$ row must reproduce the pure covariance-moderation architecture.
\item Agreement to within Monte Carlo error is the validation gate.
\end{itemize}

\subsection{4.3 Empirical power over the channel-weight grid}

\textbf{Power is monotone in both channel weights, and interior
cells are super-additive, as Section 3 predicts.}

\begin{itemize}\setlength{\itemsep}{2pt}\setlength{\parskip}{0pt}
\item Power increases strictly with each channel weight, holding the
  other fixed, across all three designs.
\item Interior cells exceed the probability-scale additive benchmark
  $p(c_A, 0) + p(0, c_B) - p(0, 0)$.
\item The design ranking (Hybrid $>$ OL+BDC $>$ CO) is preserved
  throughout the grid.
\end{itemize}

\section{5. When a Combined Architecture is Appropriate}

\textbf{The dual-channel architecture is appropriate when both a
pharmacodynamic-scaling pathway and a co-regulatory coupling pathway
are plausible for the same biomarker.}

\begin{itemize}\setlength{\itemsep}{2pt}\setlength{\parskip}{0pt}
\item Use it when independent evidence supports both channels, not
  as a default.
\item When channel weights are unknown, the $\alpha$-mixing
  parameterization spans pure A to pure B as a sensitivity analysis.
\item It carries two free parameters where one suffices, so
  calibration demands more prior information.
\end{itemize}

\section{6. Discussion}

\textbf{The combined architecture unifies the companion paper's two
architectures and explains their descriptive combined results.}

\begin{itemize}\setlength{\itemsep}{2pt}\setlength{\parskip}{0pt}
\item A and B are single-channel limits of one model, not arbitrary
  alternatives.
\item Super-additivity is a property of the power function, not an
  artifact of the calibration.
\item The analysis model is held fixed, so calibration questions are
  deferred to the companion calibration study [@pmsimstats-paper10].
\end{itemize}

\end{document}
